%───────────────────────────────────────────────────────────────────────────────
% CAPTIONS-CONFIG.TEX — Estilo de captions con paleta Catppuccin
%───────────────────────────────────────────────────────────────────────────────
% Requiere: caption, subcaption (cargados en packages.tex), tema cargado.
%───────────────────────────────────────────────────────────────────────────────

\captionsetup{
  font        = small,
  labelfont   = {bf, color=CtpBlue},         % "Figura 1.1" en bold + azul
  textfont    = {color=CtpSubtext0},          % texto del caption en gris suave
  format      = hang,                        % sangría francesa
  labelsep    = endash,                      % "Figura 1.1 –" en lugar de ":"
  skip        = 8pt,                         % espacio entre flotante y caption
  justification = justified,
}

% ─── Subfiguras ──────────────────────────────────────────────────────────────
\captionsetup[sub]{
  font        = footnotesize,
  labelfont   = {bf, color=CtpSapphire},     % "(a)" en sapphire
  textfont    = {color=CtpSubtext0},
  labelsep    = space,
  skip        = 4pt,
}

% ─── Tablas (override si se quiere diferente de figuras) ─────────────────────
\captionsetup[table]{
  position    = above,                       % caption arriba de la tabla
  labelfont   = {bf, color=CtpMauve},        % "Tabla 1.1" en mauve
  skip        = 6pt,
}

% ─── Algoritmos ──────────────────────────────────────────────────────────────
% algorithm2e maneja su propio caption; el color se configura en algorithms.tex
% con \SetAlCapSty. No se necesita nada aquí.

% ─── Listings ────────────────────────────────────────────────────────────────
\captionsetup[lstlisting]{
  labelfont   = {bf, color=CtpGreen},        % "Código 1.1" en verde
  textfont    = {color=CtpSubtext0},
  format      = hang,
  labelsep    = endash,
  skip        = 6pt,
}

% ─── Renombrar labels en español ─────────────────────────────────────────────
% "Cuadro" → "Tabla" ya se hace en tables-config.tex
% "Código fuente" para listings:
\renewcommand{\lstlistingname}{Código}
\renewcommand{\lstlistlistingname}{Índice de códigos}
